\section{Pastoral and Canonical Questions}

\subsection{Is celibacy ``only a discipline''?}

It is an ecclesiastical admission discipline, not a sacramental validity requirement or an intrinsic demand of priesthood. The word \emph{only} is misleading if it suggests theological emptiness or optional private observance. The Latin Church authoritatively judges celibacy especially fitting for presbyteral ministry and binds those ordained under the norm until competent dispensation.

\subsection{Does the existence of married priests disprove the Latin rationale?}

No. Marriage and celibacy are different gifts and signs. A married priest's sacramental character and ministerial power are complete. A celibate priest's undivided commitment can remain a genuine eschatological and pastoral sign. Each discipline needs formation suited to its real costs rather than rhetoric that demeans the other.

\subsection{May a parish expect a deacon's wife to function as an unpaid co-deacon?}

No. She consents to his ordination; she is not ordained. She retains her own vocation, employment, privacy, apostolate, limits, and marital rights. Diocesan formation may rightly include her and attend to family readiness, but coercive expectations can damage both marriage and ministry.

\subsection{Does canon 277 require a married permanent deacon and wife to abstain?}

The canon's general wording generated a credible textual argument for that conclusion. The competent Roman clarification in 2011, read with Vatican II and the papally approved 1998 diaconal norms, rejects it during the valid marriage. The scholarly dispute over formal interpretive weight should be recorded; couples should not be told that the operative Church discipline secretly demands total abstinence.

\subsection{Can a widowed permanent deacon remarry?}

Not ordinarily. Sacred orders remains a diriment impediment. Exceptional circumstances may support a petition to the Apostolic See, but no listed circumstance grants permission automatically. The deacon should seek his ordinary and a qualified canonist before entering any relationship or making civil plans.

\subsection{Are Eastern priests ``dispensed from celibacy''?}

No. A married clergy belongs to the proper and legitimate discipline of the Eastern Catholic Churches. CCEO common law honors both states and delegates admission of married men to the applicable particular law and special norms. Latin exceptions operate through a different juridic structure.

\subsection{Does a failure against chastity invalidate ministry?}

Sin does not erase ordination. It can require repentance, treatment, restrictions, removal from office, penal process, or dismissal depending on the conduct and law. Sacramental validity in a particular act has its own requirements. Abuse, coercion, and offenses involving minors or protected persons must never be reduced to private impurity.

\subsection{How should reform be discussed?}

Distinguish at least four questions: whether the Church has authority to change an admission discipline; whether change would be prudent in a place and time; what theology the existing practice expresses; and how current obligations bind persons already ordained. Evidence for married ministry answers the first but not automatically the second. The law as it is binds until changed by competent authority.

\begin{studybox}{Case-reading checklist}
Identify the person's Church \emph{sui iuris}; clerical degree; institute or ordinariate; marital history and validity; spouse's status and consent; children and dependents; whether the event is candidacy, ordination, separation, widowhood, attempted marriage, loss of state, or a penal allegation; the universal, particular, and proper law in force on the relevant date; and the authority competent to permit, dispense, investigate, or judge. Do not offer a conclusion before the facts and texts are aligned.
\end{studybox}
